\section{The object: the symplectic quadrangle \texorpdfstring{$W(3,3)$}{W(3,3)}}
\label{sec:object}

\subsection{Definition}

Let $V=\F_3^4$ with coordinates $x=(x_1,x_2,x_3,x_4)$ and the \emph{symplectic form}
\[
\langle x,y\rangle \;=\; x_1y_3+x_2y_4-x_3y_1-x_4y_2 \pmod 3 .
\]
It is bilinear, alternating ($\langle x,x\rangle=0$) and nondegenerate (only $x=0$ pairs
to zero with everything).

\begin{definition}
The \emph{points} of $\W$ are the one-dimensional subspaces of $V$ ($40$ of them). Two
distinct points $\langle x\rangle,\langle y\rangle$ are \emph{collinear} if
$\langle x,y\rangle=0$. The \emph{lines} of $\W$ are the two-dimensional subspaces $L$
on which the form vanishes identically (\emph{totally isotropic} or \emph{Lagrangian}
planes); each contains $(9-1)/2=4$ points.
\end{definition}

\begin{result}{Proposition 2.1 \hfill \tierT}
$\W$ has $40$ points and $40$ lines. Every line has $4$ points and every point lies on
$4$ lines. For a point $p$ and a line $L$ not containing it, exactly one point of $L$ is
collinear with $p$.
\end{result}

\begin{proof}
Points: $(3^4-1)/(3-1)=40$. For a point $p=\langle x\rangle$, the set
$x^\perp=\{y:\langle x,y\rangle=0\}$ is a $3$-dimensional subspace containing $x$; the
lines through $p$ are the $2$-dimensional subspaces of $x^\perp$ that contain $x$, and
these correspond to the points of $x^\perp/\langle x\rangle\cong\F_3^2$, of which there
are $(9-1)/2=4$. Double counting point--line incidences, $40\cdot4=(\#\text{lines})\cdot4$,
so there are $40$ lines. Finally, if $p\notin L$ then $x^\perp\cap L$ is a subspace of
$L$ of dimension $1$ (it cannot be all of $L$, since then $\langle x\rangle+L$ would be a
totally isotropic $3$-space, impossible in a nondegenerate $4$-space), so exactly one
point of $L$ is collinear with $p$.
\end{proof}

The last property is the defining axiom of a \emph{generalized quadrangle}: there are no
triangles made of three lines, but from any point you can always reach any line in one
step, in exactly one way. With four points per line and four lines per point, $\W$ is
the generalized quadrangle of order $(3,3)$, usually written $\mathrm{GQ}(3,3)$ \cite{PayneThas}.

\begin{figure}[htbp]
\centering
\begin{tikzpicture}[scale=0.95]
  \fill[ink] (0,0) circle (3pt) node[below left=1pt,font=\small] {$p$};
  \foreach \a/\c/\n in {90/teal/1,180/sky/2,270/gold/3,0/sage/4}{
    \draw[\c,line width=1.3pt] (0,0) -- (\a:2.6);
    \foreach \r in {0.9,1.75,2.6}{\fill[\c] (\a:\r) circle (2.4pt);} }
  \node[font=\footnotesize,align=center,text=slate] at (0,-3.3)
     {$1$ point \;+\; $4\times3=12$ neighbours \;+\; $27$ far points \;=\; $40$};
  \begin{scope}[xshift=6.2cm]
    \draw[ink!60,line width=1.1pt] (-1.8,1.4) -- (1.8,1.4) node[right,font=\small]{$L$};
    \foreach \x in {-1.5,-0.5,0.5,1.5}{\fill[ink!70] (\x,1.4) circle (2.4pt);}
    \fill[rose] (0,-1.3) circle (3pt) node[below,font=\small]{$p\notin L$};
    \draw[rose,line width=1.2pt,dashed] (0,-1.3) -- (0.5,1.4);
    \node[font=\footnotesize,text=slate,align=center] at (0,-2.5)
     {exactly one point of $L$\\is collinear with $p$};
  \end{scope}
\end{tikzpicture}
\caption{Left: the view of $\W$ from one point. Right: the quadrangle axiom.}
\label{fig:pencil}
\end{figure}

\subsection{The collinearity graph and its spectrum}

Join two points by an edge when they are collinear. Each point has $4\times3=12$
neighbours, so the graph has $40\cdot12/2=240$ edges. Two collinear points share exactly
$2$ further neighbours (the other two points of their line), and two non-collinear points
share exactly $4$ (the points of $x^\perp\cap y^\perp$, a nondegenerate plane). A graph in
which these three numbers are constant is called \emph{strongly regular}; this one is
$\SRG(40,12,2,4)$ \cite{BvM}.

\begin{result}{Proposition 2.2 \hfill \tierT}
The adjacency matrix $A$ of $\W$ has eigenvalues $12$ (once), $2$ ($24$ times) and $-4$
($15$ times).
\end{result}

\begin{proof}
For $\SRG(v,k,\lambda,\mu)$ one has $A^2=kI+\lambda A+\mu(J-I-A)$, where $J$ is the
all-ones matrix. On vectors orthogonal to the all-ones vector, $J$ acts as zero, so every
other eigenvalue $\theta$ satisfies $\theta^2=(\lambda-\mu)\theta+(k-\mu)$, i.e.
$\theta^2+2\theta-8=0$, giving $\theta\in\{2,-4\}$. Multiplicities $f,g$ satisfy
$f+g=39$ and $\tr A=12+2f-4g=0$, so $f=24$, $g=15$.
\end{proof}

Since $|-4|<2\sqrt{11}\approx6.63$, the graph is \emph{Ramanujan}: it is as good an
expander as a $12$-regular graph can be. This is why the machine of Section~10 can route
between any two of its forty nodes in two hops, and why the graph's minimum bisection is
large ($100$ of the $240$ edges, a computed value).

\begin{figure}[p]
\centering
\begin{tikzpicture}[scale=1.02,
   ff/.style={rose!55,line width=0.22pt,opacity=0.55},
   nf/.style={sky!70,line width=0.28pt,opacity=0.6},
   nn/.style={ink!70,line width=1.1pt},
   pn/.style={ink!80,line width=1.1pt}]
  % generated by make_figures.py -- W(3,3), 40 points, 240 edges
\draw[ff] (0.000,4.100) -- (4.038,0.712);
\draw[ff] (0.000,4.100) -- (2.635,3.141);
\draw[ff] (0.000,4.100) -- (3.551,-2.050);
\draw[ff] (0.000,4.100) -- (-1.402,-3.853);
\draw[ff] (0.000,4.100) -- (1.402,-3.853);
\draw[ff] (0.000,4.100) -- (-3.551,-2.050);
\draw[ff] (4.038,0.712) -- (2.635,3.141);
\draw[ff] (4.038,0.712) -- (-4.093,-0.238);
\draw[ff] (4.038,0.712) -- (-0.946,3.989);
\draw[ff] (4.038,0.712) -- (-2.982,-2.814);
\draw[ff] (2.635,3.141) -- (-2.253,-3.426);
\draw[ff] (2.635,3.141) -- (0.476,-4.072);
\draw[ff] (2.635,3.141) -- (-3.928,-1.176);
\draw[ff] (2.635,3.141) -- (-3.765,1.624);
\draw[ff] (3.551,-2.050) -- (-3.551,-2.050);
\draw[ff] (3.551,-2.050) -- (4.093,-0.238);
\draw[ff] (3.551,-2.050) -- (-2.253,-3.426);
\draw[ff] (3.551,-2.050) -- (-3.765,1.624);
\draw[ff] (3.551,-2.050) -- (-0.946,3.989);
\draw[ff] (-1.402,-3.853) -- (4.093,-0.238);
\draw[ff] (-1.402,-3.853) -- (0.476,-4.072);
\draw[ff] (-1.402,-3.853) -- (-4.093,-0.238);
\draw[ff] (1.402,-3.853) -- (4.093,-0.238);
\draw[ff] (-4.038,0.712) -- (0.000,4.100);
\draw[ff] (-4.038,0.712) -- (-1.402,-3.853);
\draw[ff] (-4.038,0.712) -- (3.289,2.448);
\draw[ff] (-4.038,0.712) -- (3.928,-1.176);
\draw[ff] (-4.038,0.712) -- (2.982,-2.814);
\draw[ff] (-4.038,0.712) -- (-3.928,-1.176);
\draw[ff] (-4.038,0.712) -- (-0.946,3.989);
\draw[ff] (-4.038,0.712) -- (-0.476,-4.072);
\draw[ff] (-2.635,3.141) -- (0.000,4.100);
\draw[ff] (-2.635,3.141) -- (1.402,-3.853);
\draw[ff] (-2.635,3.141) -- (3.928,-1.176);
\draw[ff] (-2.635,3.141) -- (-2.253,-3.426);
\draw[ff] (-2.635,3.141) -- (-3.765,1.624);
\draw[ff] (-2.635,3.141) -- (-4.093,-0.238);
\draw[ff] (-2.635,3.141) -- (2.253,-3.426);
\draw[ff] (0.946,3.989) -- (4.093,-0.238);
\draw[ff] (0.946,3.989) -- (0.476,-4.072);
\draw[ff] (0.946,3.989) -- (-0.946,3.989);
\draw[ff] (3.289,2.448) -- (-3.551,-2.050);
\draw[ff] (3.289,2.448) -- (-2.635,3.141);
\draw[ff] (3.289,2.448) -- (0.946,3.989);
\draw[ff] (3.289,2.448) -- (4.093,-0.238);
\draw[ff] (3.289,2.448) -- (2.253,-3.426);
\draw[ff] (3.289,2.448) -- (-0.476,-4.072);
\draw[ff] (3.289,2.448) -- (-2.982,-2.814);
\draw[ff] (1.840,3.664) -- (3.765,1.624);
\draw[ff] (1.840,3.664) -- (3.928,-1.176);
\draw[ff] (1.840,3.664) -- (-3.928,-1.176);
\draw[ff] (1.840,3.664) -- (-3.765,1.624);
\draw[ff] (1.840,3.664) -- (2.253,-3.426);
\draw[ff] (1.840,3.664) -- (-0.476,-4.072);
\draw[ff] (1.840,3.664) -- (-2.982,-2.814);
\draw[ff] (3.765,1.624) -- (3.551,-2.050);
\draw[ff] (3.765,1.624) -- (-1.402,-3.853);
\draw[ff] (3.765,1.624) -- (1.402,-3.853);
\draw[ff] (3.765,1.624) -- (3.928,-1.176);
\draw[ff] (3.765,1.624) -- (-4.093,-0.238);
\draw[ff] (3.765,1.624) -- (-0.946,3.989);
\draw[ff] (3.928,-1.176) -- (-3.551,-2.050);
\draw[ff] (3.928,-1.176) -- (2.982,-2.814);
\draw[ff] (3.928,-1.176) -- (-2.253,-3.426);
\draw[ff] (3.928,-1.176) -- (0.476,-4.072);
\draw[ff] (2.982,-2.814) -- (2.635,3.141);
\draw[ff] (2.982,-2.814) -- (1.402,-3.853);
\draw[ff] (2.982,-2.814) -- (0.946,3.989);
\draw[ff] (2.982,-2.814) -- (-4.093,-0.238);
\draw[ff] (2.982,-2.814) -- (-2.982,-2.814);
\draw[ff] (-2.253,-3.426) -- (0.946,3.989);
\draw[ff] (-2.253,-3.426) -- (-3.928,-1.176);
\draw[ff] (0.476,-4.072) -- (-3.551,-2.050);
\draw[ff] (0.476,-4.072) -- (-3.765,1.624);
\draw[ff] (0.476,-4.072) -- (-0.946,3.989);
\draw[ff] (-3.928,-1.176) -- (1.402,-3.853);
\draw[ff] (-3.928,-1.176) -- (4.093,-0.238);
\draw[ff] (-3.928,-1.176) -- (-0.946,3.989);
\draw[ff] (-3.765,1.624) -- (4.093,-0.238);
\draw[ff] (-1.840,3.664) -- (2.635,3.141);
\draw[ff] (-1.840,3.664) -- (-1.402,-3.853);
\draw[ff] (-1.840,3.664) -- (-3.551,-2.050);
\draw[ff] (-1.840,3.664) -- (4.093,-0.238);
\draw[ff] (-1.840,3.664) -- (1.840,3.664);
\draw[ff] (-1.840,3.664) -- (2.982,-2.814);
\draw[ff] (-1.840,3.664) -- (-3.289,2.448);
\draw[ff] (-1.840,3.664) -- (-0.476,-4.072);
\draw[ff] (-4.093,-0.238) -- (0.946,3.989);
\draw[ff] (-4.093,-0.238) -- (-3.765,1.624);
\draw[ff] (-3.289,2.448) -- (4.038,0.712);
\draw[ff] (-3.289,2.448) -- (1.402,-3.853);
\draw[ff] (-3.289,2.448) -- (-3.551,-2.050);
\draw[ff] (-3.289,2.448) -- (0.946,3.989);
\draw[ff] (-3.289,2.448) -- (3.765,1.624);
\draw[ff] (-3.289,2.448) -- (-2.253,-3.426);
\draw[ff] (-3.289,2.448) -- (2.253,-3.426);
\draw[ff] (2.253,-3.426) -- (4.038,0.712);
\draw[ff] (2.253,-3.426) -- (-1.402,-3.853);
\draw[ff] (2.253,-3.426) -- (0.476,-4.072);
\draw[ff] (2.253,-3.426) -- (-3.928,-1.176);
\draw[ff] (-0.476,-4.072) -- (4.038,0.712);
\draw[ff] (-0.476,-4.072) -- (3.551,-2.050);
\draw[ff] (-0.476,-4.072) -- (-2.253,-3.426);
\draw[ff] (-0.476,-4.072) -- (-4.093,-0.238);
\draw[ff] (-2.982,-2.814) -- (1.402,-3.853);
\draw[ff] (-2.982,-2.814) -- (-3.551,-2.050);
\draw[ff] (-2.982,-2.814) -- (-3.765,1.624);
\draw[ff] (-2.982,-2.814) -- (-0.946,3.989);
\draw[nf] (0.375,1.504) -- (0.000,4.100);
\draw[nf] (0.375,1.504) -- (4.038,0.712);
\draw[nf] (0.375,1.504) -- (2.635,3.141);
\draw[nf] (0.375,1.504) -- (0.946,3.989);
\draw[nf] (0.375,1.504) -- (4.093,-0.238);
\draw[nf] (0.375,1.504) -- (3.289,2.448);
\draw[nf] (0.375,1.504) -- (1.840,3.664);
\draw[nf] (0.375,1.504) -- (3.765,1.624);
\draw[nf] (0.375,1.504) -- (3.928,-1.176);
\draw[nf] (0.000,4.100) -- (1.504,-0.375);
\draw[nf] (4.038,0.712) -- (1.504,0.375);
\draw[nf] (4.038,0.712) -- (0.375,-1.504);
\draw[nf] (-1.504,0.375) -- (0.000,4.100);
\draw[nf] (-1.504,0.375) -- (3.551,-2.050);
\draw[nf] (-1.504,0.375) -- (-3.551,-2.050);
\draw[nf] (-1.504,0.375) -- (0.946,3.989);
\draw[nf] (-1.504,0.375) -- (1.840,3.664);
\draw[nf] (-1.504,0.375) -- (2.982,-2.814);
\draw[nf] (-1.504,0.375) -- (-3.928,-1.176);
\draw[nf] (-1.504,0.375) -- (-4.093,-0.238);
\draw[nf] (-1.504,0.375) -- (2.253,-3.426);
\draw[nf] (-0.375,-1.504) -- (0.000,4.100);
\draw[nf] (-0.375,-1.504) -- (-1.402,-3.853);
\draw[nf] (-0.375,-1.504) -- (0.946,3.989);
\draw[nf] (-0.375,-1.504) -- (-2.253,-3.426);
\draw[nf] (-0.375,-1.504) -- (-3.765,1.624);
\draw[nf] (-0.375,-1.504) -- (-2.982,-2.814);
\draw[nf] (1.504,-0.375) -- (1.402,-3.853);
\draw[nf] (1.504,-0.375) -- (0.946,3.989);
\draw[nf] (1.504,-0.375) -- (0.476,-4.072);
\draw[nf] (1.504,-0.375) -- (-0.946,3.989);
\draw[nf] (-4.038,0.712) -- (-0.375,-1.504);
\draw[nf] (-4.038,0.712) -- (-0.375,1.504);
\draw[nf] (-4.038,0.712) -- (-1.550,0.000);
\draw[nf] (-4.038,0.712) -- (1.504,0.375);
\draw[nf] (-2.635,3.141) -- (1.504,-0.375);
\draw[nf] (-2.635,3.141) -- (-0.375,1.504);
\draw[nf] (-2.635,3.141) -- (-0.000,-1.550);
\draw[nf] (-2.635,3.141) -- (-1.504,-0.375);
\draw[nf] (0.000,1.550) -- (3.551,-2.050);
\draw[nf] (0.000,1.550) -- (-1.402,-3.853);
\draw[nf] (0.000,1.550) -- (1.402,-3.853);
\draw[nf] (0.000,1.550) -- (2.982,-2.814);
\draw[nf] (0.000,1.550) -- (-2.253,-3.426);
\draw[nf] (0.000,1.550) -- (0.476,-4.072);
\draw[nf] (0.000,1.550) -- (-2.982,-2.814);
\draw[nf] (-0.375,1.504) -- (-3.551,-2.050);
\draw[nf] (-0.375,1.504) -- (-3.928,-1.176);
\draw[nf] (-0.375,1.504) -- (-3.765,1.624);
\draw[nf] (-0.375,1.504) -- (-4.093,-0.238);
\draw[nf] (-0.375,1.504) -- (-0.946,3.989);
\draw[nf] (3.289,2.448) -- (-1.550,0.000);
\draw[nf] (3.289,2.448) -- (-0.000,-1.550);
\draw[nf] (3.289,2.448) -- (1.550,0.000);
\draw[nf] (1.840,3.664) -- (-0.375,-1.504);
\draw[nf] (1.840,3.664) -- (1.504,-0.375);
\draw[nf] (3.765,1.624) -- (-0.000,-1.550);
\draw[nf] (3.765,1.624) -- (1.550,0.000);
\draw[nf] (3.928,-1.176) -- (1.504,0.375);
\draw[nf] (3.928,-1.176) -- (0.375,-1.504);
\draw[nf] (-1.550,0.000) -- (2.635,3.141);
\draw[nf] (-1.550,0.000) -- (1.402,-3.853);
\draw[nf] (-1.550,0.000) -- (3.765,1.624);
\draw[nf] (-1.550,0.000) -- (0.476,-4.072);
\draw[nf] (-1.550,0.000) -- (-3.765,1.624);
\draw[nf] (-1.550,0.000) -- (-3.289,2.448);
\draw[nf] (-1.550,0.000) -- (-0.476,-4.072);
\draw[nf] (-0.000,-1.550) -- (2.635,3.141);
\draw[nf] (-0.000,-1.550) -- (3.551,-2.050);
\draw[nf] (-0.000,-1.550) -- (2.982,-2.814);
\draw[nf] (-0.000,-1.550) -- (-0.946,3.989);
\draw[nf] (1.550,0.000) -- (2.635,3.141);
\draw[nf] (1.550,0.000) -- (-1.402,-3.853);
\draw[nf] (1.550,0.000) -- (-3.551,-2.050);
\draw[nf] (1.550,0.000) -- (-2.253,-3.426);
\draw[nf] (1.550,0.000) -- (-3.928,-1.176);
\draw[nf] (1.550,0.000) -- (-4.093,-0.238);
\draw[nf] (1.550,0.000) -- (-2.982,-2.814);
\draw[nf] (2.982,-2.814) -- (1.504,0.375);
\draw[nf] (-1.840,3.664) -- (1.504,-0.375);
\draw[nf] (-1.840,3.664) -- (-0.375,1.504);
\draw[nf] (-1.840,3.664) -- (-0.000,-1.550);
\draw[nf] (-1.840,3.664) -- (-1.504,-0.375);
\draw[nf] (-1.504,-0.375) -- (4.038,0.712);
\draw[nf] (-1.504,-0.375) -- (-1.402,-3.853);
\draw[nf] (-1.504,-0.375) -- (4.093,-0.238);
\draw[nf] (-1.504,-0.375) -- (3.928,-1.176);
\draw[nf] (-1.504,-0.375) -- (-2.253,-3.426);
\draw[nf] (-1.504,-0.375) -- (-0.946,3.989);
\draw[nf] (-1.504,-0.375) -- (-2.982,-2.814);
\draw[nf] (1.504,0.375) -- (3.551,-2.050);
\draw[nf] (1.504,0.375) -- (4.093,-0.238);
\draw[nf] (1.504,0.375) -- (-3.765,1.624);
\draw[nf] (0.375,-1.504) -- (1.402,-3.853);
\draw[nf] (0.375,-1.504) -- (-3.551,-2.050);
\draw[nf] (0.375,-1.504) -- (4.093,-0.238);
\draw[nf] (0.375,-1.504) -- (0.476,-4.072);
\draw[nf] (0.375,-1.504) -- (-3.928,-1.176);
\draw[nf] (0.375,-1.504) -- (-4.093,-0.238);
\draw[nf] (-3.289,2.448) -- (-0.375,-1.504);
\draw[nf] (-3.289,2.448) -- (-0.375,1.504);
\draw[nf] (-3.289,2.448) -- (1.504,0.375);
\draw[nf] (2.253,-3.426) -- (0.000,1.550);
\draw[nf] (2.253,-3.426) -- (-0.000,-1.550);
\draw[nf] (2.253,-3.426) -- (1.504,0.375);
\draw[nf] (-0.476,-4.072) -- (1.504,-0.375);
\draw[nf] (-0.476,-4.072) -- (0.000,1.550);
\draw[nf] (-0.476,-4.072) -- (0.375,-1.504);
\draw[nn] (0.375,1.504) -- (0.000,1.550);
\draw[nn] (0.375,1.504) -- (-0.375,1.504);
\draw[nn] (-1.504,0.375) -- (-1.504,-0.375);
\draw[nn] (-0.375,-1.504) -- (-0.000,-1.550);
\draw[nn] (-0.375,-1.504) -- (0.375,-1.504);
\draw[nn] (1.504,-0.375) -- (1.550,0.000);
\draw[nn] (1.504,-0.375) -- (1.504,0.375);
\draw[nn] (0.000,1.550) -- (-0.375,1.504);
\draw[nn] (-1.550,0.000) -- (-1.504,0.375);
\draw[nn] (-1.550,0.000) -- (-1.504,-0.375);
\draw[nn] (-0.000,-1.550) -- (0.375,-1.504);
\draw[nn] (1.550,0.000) -- (1.504,0.375);
\draw[pn] (0.375,1.504) -- (0.000,0.000);
\draw[pn] (-1.504,0.375) -- (0.000,0.000);
\draw[pn] (-0.375,-1.504) -- (0.000,0.000);
\draw[pn] (0.000,0.000) -- (1.504,-0.375);
\draw[pn] (0.000,0.000) -- (0.000,1.550);
\draw[pn] (0.000,0.000) -- (-0.375,1.504);
\draw[pn] (0.000,0.000) -- (-0.000,-1.550);
\draw[pn] (0.000,0.000) -- (1.550,0.000);
\draw[pn] (0.000,0.000) -- (-1.504,-0.375);
\draw[pn] (0.000,0.000) -- (1.504,0.375);
\draw[pn] (0.000,0.000) -- (0.375,-1.504);
\draw[pn] (-1.550,0.000) -- (0.000,0.000);
\fill[teal] (0.375,1.504) circle (2.6pt);
\fill[teal] (0.000,1.550) circle (2.6pt);
\fill[teal] (-0.375,1.504) circle (2.6pt);
\fill[sky] (-1.504,0.375) circle (2.6pt);
\fill[sky] (-1.550,0.000) circle (2.6pt);
\fill[sky] (-1.504,-0.375) circle (2.6pt);
\fill[gold] (-0.375,-1.504) circle (2.6pt);
\fill[gold] (-0.000,-1.550) circle (2.6pt);
\fill[gold] (0.375,-1.504) circle (2.6pt);
\fill[sage] (1.504,-0.375) circle (2.6pt);
\fill[sage] (1.550,0.000) circle (2.6pt);
\fill[sage] (1.504,0.375) circle (2.6pt);
\fill[rose] (0.000,4.100) circle (2.2pt);
\fill[rose] (0.946,3.989) circle (2.2pt);
\fill[rose] (1.840,3.664) circle (2.2pt);
\fill[rose] (2.635,3.141) circle (2.2pt);
\fill[rose] (3.289,2.448) circle (2.2pt);
\fill[rose] (3.765,1.624) circle (2.2pt);
\fill[rose] (4.038,0.712) circle (2.2pt);
\fill[rose] (4.093,-0.238) circle (2.2pt);
\fill[rose] (3.928,-1.176) circle (2.2pt);
\fill[rose] (3.551,-2.050) circle (2.2pt);
\fill[rose] (2.982,-2.814) circle (2.2pt);
\fill[rose] (2.253,-3.426) circle (2.2pt);
\fill[rose] (1.402,-3.853) circle (2.2pt);
\fill[rose] (0.476,-4.072) circle (2.2pt);
\fill[rose] (-0.476,-4.072) circle (2.2pt);
\fill[rose] (-1.402,-3.853) circle (2.2pt);
\fill[rose] (-2.253,-3.426) circle (2.2pt);
\fill[rose] (-2.982,-2.814) circle (2.2pt);
\fill[rose] (-3.551,-2.050) circle (2.2pt);
\fill[rose] (-3.928,-1.176) circle (2.2pt);
\fill[rose] (-4.093,-0.238) circle (2.2pt);
\fill[rose] (-4.038,0.712) circle (2.2pt);
\fill[rose] (-3.765,1.624) circle (2.2pt);
\fill[rose] (-3.289,2.448) circle (2.2pt);
\fill[rose] (-2.635,3.141) circle (2.2pt);
\fill[rose] (-1.840,3.664) circle (2.2pt);
\fill[rose] (-0.946,3.989) circle (2.2pt);
\fill[ink] (0,0) circle (3.4pt);

  \node[font=\sffamily\small\bfseries,text=ink] at (0,-4.85) {$W(3,3)$: $40$ points, all $240$ edges};
\end{tikzpicture}

\medskip
\begin{tikzpicture}
  \foreach \c/\t/\x in {ink/{$p$ and its $4$ lines ($12+12$ edges)}/0,
                        sky/{neighbour to far point ($108$)}/5.4,
                        rose/{far to far ($108$)}/10.1}{
    \draw[\c,line width=1.4pt] (\x,0) -- (\x+0.5,0);
    \node[anchor=west,font=\sffamily\scriptsize] at (\x+0.6,0) {\t};}
\end{tikzpicture}
\caption{A portrait of $\W$, computed from the symplectic form and drawn from the point
of view of one point $p$ (centre). Its $12$ neighbours lie on four lines (four colours,
three points each); the $27$ far points sit on the outer ring, ordered by which neighbour
they see on each of the four lines --- by the quadrangle axiom every far point sees exactly
one. The edge counts $12+12+108+108=240$ are forced by $\SRG(40,12,2,4)$.}
\label{fig:portrait}
\end{figure}

\begin{plainwords}
``Strongly regular'' means that the geometry looks exactly the same from every point, and
from every pair of points. There is no privileged measurement. The three eigenvalues are
the geometry's ``vibration frequencies''; that there are only three is the algebraic
signature of this perfect evenness.
\end{plainwords}

\subsection{What the points and lines are physically}

Let $W(x)=X^{x_1}Z^{x_3}\otimes X^{x_2}Z^{x_4}$ be the two-qutrit Pauli operator with
label $x\in\F_3^4$. A direct calculation from $ZX=\omega XZ$ gives
\begin{equation}\label{eq:commutation}
W(x)\,W(y)\;=\;\omega^{\langle x,y\rangle}\,W(y)\,W(x).
\end{equation}
So the symplectic form is not an arbitrary choice: it \emph{is} the commutation phase of
two-qutrit measurements.

\begin{physical}[What it means physically: contexts]
\begin{itemize}[leftmargin=1.2em]
\item A \textbf{point} is a measurement direction: the cyclic group generated by one
      Pauli operator, which is the same whether it is run forwards or backwards.
\item Two points are \textbf{collinear} when their operators commute and can be measured
      at the same time without disturbing each other.
\item A \textbf{line} is a maximal set of mutually compatible measurements: an abelian
      group of $9$ Pauli operators whose common eigenvectors form a basis of the
      two-qutrit space $\C^9$. Physicists call this a \emph{stabilizer group}, and each
      of its $9$ joint eigenstates a \emph{stabilizer state}. There are $40$ lines and
      so $40\times9=360$ stabilizer states --- the standard count for two qutrits.
\item The \textbf{quadrangle axiom} says: given any measurement and any complete
      measurement context not containing it, exactly one member of the context is
      compatible with it.
\end{itemize}
\end{physical}

The quantum operations that map Pauli operators to Pauli operators (up to phase) form the
two-qutrit \emph{Clifford group} \cite{GHW}. Modulo the Paulis themselves and phases, it acts on the
labels $x$ by exactly the linear maps preserving $\langle\,,\rangle$: the symplectic
group $\Sp(4,3)$. So the symmetries of the geometry are, literally, the ``cheap'' quantum
gates.

\subsection{Two groups of order 51,840, which are not the same group}

The symmetry groups that appear around $\W$ are easy to confuse, and confusing them has
cost this programme real work (Section~\ref{sec:e8-trap}). Adapting the blueprint's
account:
\[
\PSp(4,3)\cong U_4(2)\ (\text{order }25\,920),\qquad
\Sp(4,3)\cong 2.U_4(2),\qquad
\PGSp(4,3)\cong U_4(2){:}2\cong W(E_6).
\]
The last two both have order $51\,840$ but are \emph{not isomorphic}: $\Sp(4,3)$ has a
centre of order two (the map $x\mapsto -x$), while $\PGSp(4,3)$ --- the full automorphism
group of $\W$ --- has trivial centre and is the Weyl group of the exceptional root system
$E_6$. \tierT

\begin{plainwords}
The two groups have the same size but different ``shapes''. One of them contains a
special element that commutes with everything --- a global sign flip. Whether that sign
is visible or invisible turns out to decide whether the structure can tell left from
right. This is why the same forty points can \emph{host} a handedness without being able
to \emph{choose} one (Section~\ref{sec:chirality}).
\end{plainwords}

\subsection{The dual picture and a firewall}

Swap the roles of points and lines: the $40$ lines of $\W$, with two lines ``adjacent''
when they meet, form the dual quadrangle, which is the parabolic quadric $Q(4,3)$. For odd
$q$ the quadrangle $W(3,q)$ is \emph{not} isomorphic to its dual, although both have
collinearity graph parameters $\SRG(40,12,2,4)$. This distinction matters several times
below: the finite space-time of Section~4 lives on the \emph{line} side, while the
qutrit operators live on the \emph{point} side. Repeated coincidences of counts between the
two sides (for instance $27$ vertices, $108$ edges, $36$ triangles) do not make them the
same object --- the repository calls this the \emph{point/line firewall}, and it has
caught several tempting mistakes. \tierC\ (\cert{w33\_20260924\_history\_pointline\_cospectral\_firewall.py})

\subsection{Why three?}

A sceptical reader should ask whether ``three'' was chosen because it works. The
repository records five equations in which $q$ appears as an unknown integer and each has
exactly one positive solution, $q=3$. Two can be checked by hand:
\begin{itemize}
\item The number of edges of the collinearity graph of $W(3,q)$ is
      $q(q+1)^2(q^2+1)/2$; the number of elements of $\F_{q^5}$ outside $\F_q$ is $q^5-q$.
      Setting them equal gives $2(q-1)=q+1$, so $q=3$: both equal $240$.
\item The number of points not collinear with a given point is $q^3$; it equals $27$,
      the dimension of the fundamental representation of $E_6$, only for $q=3$.
\end{itemize}

\begin{boundary}
These are coincidences of integers, not derivations. They do not prove that three is
\emph{necessary} for anything physical. They do show that the programme's most unusual
choice --- three-level systems rather than two-level ones --- is the one value at which
several independent structures meet. \tierB
\end{boundary}
